\scanpage{35}
\mspage{9/}

définit donc (avec la paire $\{z,z'\}$) une métrique sphérique, ce qui sera OK.

On fera attention \inserted{cependant} que $\sigma$ n'est pas unique, il y a une ambiguïté d'ordre $2$, on peut le remplacer par $\sigma'=\varepsilon\cdot\sigma$, où $\varepsilon$ est l'élément d'ordre $2$ de $T$. On trouve en fait que $\{\sigma,\sigma'\}$ est formé d'une inversion et de un antipodisme.

\uline{Dém} OPS $X=$ \deleted{\unsure{S|P}}sphère de Riemann, $\{z,z'\}=\{\infty,0\}$, alors
\[
T=\{z\mapsto az\mid a\in\mathbb U\ \text{i.e.}\ a\in\mathbb C,\ |a|=1\ \text{i.e.}\ a\bar a=1\}
\]
Les \deleted{\unsure{anti-au}} autom. conformes anti-hol. de $\mathbb C^*$ sont
\[
\begin{aligned}
z&\mapsto a\bar z\quad(a\in\mathbb C^*)&&\text{ceux qui fixent }0,\infty\\
z&\mapsto b/\bar z\quad(b\in\mathbb C^*)&&\rule{2cm}{.35pt}\ \text{intervertissent }0,\infty
\end{aligned}
\]
On voit \deleted{\unsure{donc|tout}} de suite que ce sont les \uline{seconds} qui commutent à $T$. \deleted{Une telle} Si
\[
\sigma(z)=b/\bar z,\quad\text{on a}\quad \sigma^2(z)=b/\bar b\,z,
\]
donc $\sigma$ involutif \unsure{dss|sss|ssi} $b/\bar b=1$ i.e. $b\in\mathbb R^*$. Les pts fixes de $\sigma$ sont donnés par
\[
z\bar z=b
\]
\deleted{\unsure{C'est}} donc $\sigma$ est une inversion resp. un antipodisme, suivant que \deleted{$b\in\mathbb R$} $b>0$ ou $b<0$.

Si $\sigma(z)=b/\bar z$, on aura $\sigma'(z)=-b/\bar z$, \deleted{\unsure{donc|de}} donc si \deleted{\unsure{l'une}} \inserted{$\sigma$} est une inversion, $\sigma'$ est un antipodisme, et vice versa. OPS les notations choisies de telle façon que $b>0$, i.e. $\sigma$ inversion, donc $\sigma$ définit une métrique sphérique ayant pour grand cercle défini par $\unsure{\mathscr C|\mathscr E|C}=\{z\in\mathbb C\mid z\bar z=b\}$.

% archive sheet number: 37
